% ===========================================================
%  CHAPTER V — IMPLEMENTATION AND EXPERIMENTAL METHODOLOGY
% ===========================================================
\chapter{Implementation and Experimental Methodology}
\label{chap:impl}

% -----------------------------------------------------------
\section{Technology Stack}

\begin{table}[H]
  \centering
  \caption{Software actually used, at the versions the reported results were
           produced with.}
  \label{tab:stack}
  \rowcolors{2}{lightblue!30}{white}
  \small
  \setlength{\tabcolsep}{5pt}
  \begin{tabularx}{\textwidth}{@{} L{3.5cm} C{1.9cm} X @{}}
    \toprule
    \rowcolor{lightblue}
    \textbf{Library} & \textbf{Version} & \textbf{Usage} \\
    \midrule
    Python             & 3.12.10 & \\
    MuJoCo             & 3.7.0   & Physics simulation engine~\parencite{Todorov2012} \\
    Gymnasium          & 1.2.3   & Standard \RL{} environment interface \\
    Stable-Baselines3  & 2.8.0   & DDPG/TD3/SAC/PPO implementation~\parencite{Raffin2021} \\
    PyTorch            & 2.11.0+cu126 & Deep-learning backend. The CUDA build is installed, but \texttt{config.DEVICE} defaults to \texttt{cpu} and every reported result was produced on CPU; the GPU path was used only to schedule independent jobs concurrently while producing the runs. \\
    NumPy              & 2.2.6   & Numerical computation \\
    SciPy              & 1.13.0  & Static-optimisation QP (SLSQP), NNLS \\
    scikit-learn       & 1.4.2   & NMF decomposition (synergies) \\
    Optuna             & 4.9.0   & Bayesian hyperparameter optimisation~\parencite{Akiba2019} \\
    Matplotlib         & 3.8.4   & Figures \\
    \bottomrule
  \end{tabularx}
\end{table}

The Stable-Baselines3 version is material rather than incidental. The replay-ratio
behaviour documented in \cref{sec:replaybug} --- that one rollout iteration over
an $n$-environment vector is followed by \texttt{gradient\_steps} updates, not
$n$ of them --- is a property of this implementation, and the defect it caused was
identified by measuring that library's behaviour directly rather than by reading
its documentation.

TensorBoard logging is deliberately disabled: its asynchronous event-file writer
raced with the cloud-synchronisation client on the working directory and crashed
training. Loss histories are read from the in-memory logger instead, so no
diagnostic information is lost.

% -----------------------------------------------------------
\section{Hyperparameters}

\subsection{Optimisation Procedure}

A single Optuna study was run, on \SAC only. It used TPE sampling with a
\textit{MedianPruner} over \textbf{16 trials of $10^5$ steps} each, on a training
seed fixed at 0 so that trials are comparable, with sampler seeds differing per
worker. Thirteen trials completed and three were pruned. Trials were scored by
mean final reaching error over 20 deterministic episodes on an evaluation
environment with domain randomisation disabled --- that is, under the protocol the
thesis reports, not the one training uses.

Three properties of this study are worth stating because they differ from what a
reader might assume.

\textbf{The search space excludes the network architecture.} The reported policy
footprint --- 393\,750 parameters, and the inference latency of
\cref{sec:compute_efficiency} --- is measured for the $[256, 256]$ architecture.
Searching the architecture would have invalidated those figures silently.

\textbf{The replay ratio is fixed, not searched.} Expressing it as a raw
\texttt{gradient\_steps} value would have allowed the sampler to re-enter the
under-training regime that \cref{sec:replaybug} identifies as a defect, and would
have made trials incomparable with the pilot runs. It is held at twice the
environment count throughout.

\textbf{The reward weights were not searched.} They are hand-set
(\cref{sec:reward}); only learning hyperparameters were optimised. Earlier
versions of this work described the reward weights as Optuna-optimised over 100
trials, which did not occur.

\begin{table}[H]
  \centering
  \caption{Optuna search space. \SAC only; 16 trials of $10^5$ steps.}
  \label{tab:optuna_space}
  \rowcolors{2}{lightblue!30}{white}
  \small
  \begin{tabularx}{\textwidth}{@{} L{4.2cm} L{4.2cm} X @{}}
    \toprule
    \rowcolor{lightblue}
    \textbf{Hyperparameter} & \textbf{Domain} & \textbf{Distribution} \\
    \midrule
    Learning rate        & $[5\times10^{-5},\, 10^{-3}]$   & Log-uniform \\
    $\gamma$             & $[0.95,\, 0.999]$               & Uniform \\
    Batch size           & $\{128, 256, 512\}$             & Categorical \\
    $\tau$ (Polyak)      & $[10^{-3},\, 2\times 10^{-2}]$  & Log-uniform \\
    Target entropy       & $[-27,\, -4.5]$                 & Uniform \\
    \bottomrule
  \end{tabularx}
\end{table}

\begin{table}[H]
  \centering
  \caption{Hyperparameters used. Only the \SAC ``tuned'' column was optimised;
           every other configuration uses library defaults. This asymmetry is a
           limitation of the comparison, discussed in \cref{ch:conclusion}.}
  \label{tab:hparams}
  \rowcolors{2}{lightblue!30}{white}
  \small
  \begin{tabular}{@{} L{3.6cm} C{2.2cm} C{2.2cm} C{1.8cm} C{1.8cm} C{1.8cm} @{}}
    \toprule
    \rowcolor{lightblue}
    \textbf{Hyperparameter} & \textbf{SAC tuned} & \textbf{SAC def.} & \textbf{TD3} & \textbf{DDPG} & \textbf{PPO} \\
    \midrule
    Learning rate      & $4.40\times10^{-4}$ & $3\times10^{-4}$ & $3\times10^{-4}$ & $3\times10^{-4}$ & $3\times10^{-4}$ \\
    Batch size         & 128       & 512       & 256    & 256    & 64 \\
    Buffer size        & $3\times10^5$ & $3\times10^5$ & $3\times10^5$ & $3\times10^5$ & --- \\
    $\gamma$           & 0.957     & 0.99      & 0.99   & 0.99   & 0.99 \\
    $\tau$ (Polyak)    & 0.0188    & 0.005     & 0.005  & 0.005  & --- \\
    Target entropy     & $-19.6$   & auto ($-9$) & ---  & ---    & --- \\
    Gradient steps     & 8         & 8         & 8      & 8      & --- \\
    Learning starts    & 4\,000    & 4\,000    & 4\,000 & 4\,000 & --- \\
    Action noise $\sigma$ & ---    & ---       & 0.1    & 0.1    & --- \\
    Policy delay       & ---       & ---       & 2      & ---    & --- \\
    n\_steps           & ---       & ---       & ---    & ---    & 2048 \\
    n\_epochs          & ---       & ---       & ---    & ---    & 10 \\
    $\lambda$ (GAE)    & ---       & ---       & ---    & ---    & 0.95 \\
    Clip $\varepsilon$ & ---       & ---       & ---    & ---    & 0.20 \\
    \midrule
    Network            & \multicolumn{5}{c}{$[256, 256]$, shared across all configurations} \\
    \bottomrule
  \end{tabular}
\end{table}

The gradient-steps value of 8 corresponds to a replay ratio of approximately
1.87 updates per environment step, given four parallel environments. It is
applied uniformly to every off-policy configuration, so it is not a confound
between them; \PPO is on-policy and has no equivalent parameter.

% -----------------------------------------------------------
\section{Experimental Protocol}

\subsection{Training and Evaluation}
\label{sec:compute}

Two training budgets are used. The algorithm comparison of
\cref{sec:benchmark} trains five configurations for $10^5$ steps with
\textbf{three seeds} each ($\{0,1,2\}$). The best configuration is additionally
retrained from initialisation for $3 \times 10^5$ steps. That configuration was
subsequently replicated across \textbf{five seeds} (\cref{sec:replication}), and
it is those runs that the force analysis of \cref{ch:forces} examines; the
per-muscle breakdown of \cref{tab:permuscle} is from seed~0 alone, but every
aggregate the chapter argues from is a five-seed mean.

The seed count is a compute constraint, not a methodological choice, and it is
the principal limitation of this work.
\textcite{Henderson2018} and \textcite{Agarwal2021} recommend ten seeds or more
precisely because estimation variance at three is large enough to make algorithm
comparisons inconclusive. That recommendation is not met here. The consequence is
stated explicitly wherever a comparison is reported: no significance test is
performed anywhere in this thesis, and differences smaller than the observed
seed-to-seed spread are reported as not supporting an ordering. Reaching ten
seeds at the corrected replay ratio (\cref{sec:replaybug}) would require
approximately 148~hours on the available hardware.

During training a learning curve is recorded every $100\,000$ steps
(	exttt{config.EVAL\_FREQ}) over five deterministic episodes. This is a monitoring signal only; it is
deliberately cheap, it is noisy at that episode count, and no value from it is
reported as a result.

Final evaluation uses \textbf{50 deterministic episodes} with domain
randomisation disabled, on an environment seeded independently of the one used
for best-checkpoint selection during training, so that the reported figure is not
the one selection was performed on.

The metrics actually computed are those defined in \cref{tab:metrics}: final
error, closest approach, drift, the graded proximity rates, energy, blow-up rate
and episode length; the co-contraction index of \cref{sec:icc}; the muscle-force
agreement statistics of \cref{ch:forces}; and inference latency. Perturbation
robustness, muscle jerk and the domain-randomisation breakdown described in
earlier versions of this work were not measured and are not reported.

\subsection{Statistical Treatment}
\label{sec:stats_protocol}

\textbf{No inferential statistics are reported in this thesis.} With three seeds
the two-sided Wilcoxon signed-rank test has a minimum attainable $p$-value of
$0.25$, which cannot reach any conventional significance threshold; computing one
would convey false precision. Earlier versions of this work reported Wilcoxon
tests with Bonferroni correction, rank-biserial effect sizes and BCa bootstrap
intervals over $n=10$ seeds. None of those analyses was performed, and they are
withdrawn (\cref{ch:conclusion}).

What is reported instead is the mean and standard deviation across seeds, stated
together, with the explicit convention that a difference comparable to or smaller
than the spread is treated as unresolved. Where a single-seed result is given ---
as for the force analysis --- it is identified as such at the point of use.

For the force comparison specifically, one distributional control \emph{is}
computed, because it is required for the statistic to be interpretable at all:
the cosine similarity between non-negative force vectors has a floor well above
zero, so a permutation null is constructed by shuffling, at each timestep, which
muscle each force magnitude belongs to (\cref{tab:forceoverall}).


\subsection{Defects in the evaluation path, and what they changed}
\label{sec:evaldefects}

A late audit of the evaluation code found four defects. None produced an
arithmetically wrong number, but two changed which episodes the reported figures
describe, so every evaluation in this work was repeated with the corrected code.

\begin{enumerate}[leftmargin=1.5em, itemsep=3pt]
\item \textbf{The evaluation seed was dead code.} \texttt{evaluate(seed=\ldots)}
never referenced its own argument, and \texttt{load\_model} fixed the environment
seed at zero. Verified: two loads with different seed arguments return
byte-identical results.

\item \textbf{\texttt{reset(seed=\ldots)} does not reseed target sampling.}
Targets are drawn from a generator seeded only in \texttt{\_\_init\_\_};
Gymnasium's \texttt{super().reset(seed=\ldots)} seeds a different generator that
this environment never uses. The target sequence can therefore only be changed at
construction.

\item \textbf{Two resets per episode.} The evaluation loop reset the vectorised
environment at the top of each iteration, while \texttt{DummyVecEnv} also
auto-resets on termination. Measured: the inner environment was reset twice per
episode, so evaluation scored on every \emph{other} target --- a different test
set from the one the force and conventional-controller analyses used, which made
those comparisons unpaired.

\item \textbf{A diverged episode could be scored a success.} On divergence the
joint speed is forced to zero, which automatically satisfies the ``settled'' half
of the success criterion. An episode blowing up within \SI{2}{\centi\meter} of
the target would have been recorded as a success and paid the success bonus.
Latent --- the divergence rate is zero throughout --- but a reward exploit of
precisely the kind \cref{sec:exploit} documents.
\end{enumerate}

\subsection{Uncertainty was understated: two components, not one}
\label{sec:variance}

The first three defects share a consequence. Every configuration was evaluated on
one fixed draw of 50 targets, so every $\pm$ previously reported was
\textbf{training-seed spread on a single test set}. The contribution of the
target draw itself was never measured.

It is not negligible. Evaluating one policy across five target seeds gives a
final error of $0.1130 \pm 0.0146$\,m, against $0.1033 \pm 0.0034$\,m across five
\emph{training} seeds on a fixed test set --- roughly four times the spread. The
graded rates move further still: \texttt{reach\_5cm} ranges from $0.16$ to $0.34$
across target draws, and the value previously reported was the largest of the
five.

Every policy was therefore re-evaluated across three target seeds, and the two
components are now reported separately.

\begin{table}[H]
\centering
\caption{Variance decomposition after re-evaluation. $\sigma_{\text{train}}$ is
the spread across training seeds with targets averaged out;
$\sigma_{\text{target}}$ the spread across target draws with training seeds
averaged out.}
\small
\begin{tabular}{@{} L{2.4cm} L{3.4cm} C{2.0cm} C{2.4cm} C{2.4cm} C{2.0cm} @{}}
\toprule
\textbf{Config.} & \textbf{Metric} & \textbf{mean} & $\boldsymbol{\sigma_{\text{train}}}$ & $\boldsymbol{\sigma_{\text{target}}}$ & \textbf{total} \\
\midrule
$w_2 = 1$  & final error & 0.1074 & 0.0043 & 0.0066 & 0.0079 \\
           & closest     & 0.0734 & 0.0050 & 0.0031 & 0.0058 \\
           & ends $<$5\,cm & 0.260 & 0.0374 & 0.0666 & 0.0764 \\
\midrule
$w_2 = 5$  & final error & 0.1114 & 0.0040 & 0.0035 & 0.0053 \\
           & closest     & 0.0750 & 0.0037 & 0.0034 & 0.0050 \\
\midrule
$w_2 = 20$ & final error & 0.1161 & 0.0095 & 0.0053 & 0.0108 \\
           & closest     & 0.0737 & 0.0073 & 0.0015 & 0.0074 \\
           & ends $<$5\,cm & 0.192 & 0.0477 & 0.0174 & 0.0508 \\
\bottomrule
\end{tabular}
\label{tab:variance}
\end{table}

For most quantities $\sigma_{\text{target}}$ is comparable to or larger than
$\sigma_{\text{train}}$. \textbf{Which targets a policy is tested on matters at
least as much as which seed it was trained from}, and reporting only the latter
understated the uncertainty on every absolute figure in this work.

\subsection*{What this does \emph{not} undermine}

One consequence of the defect is favourable. Because the environment seed was
fixed, every configuration was evaluated on \emph{identical} targets, so all
between-configuration comparisons are paired. Paired comparison removes the
target draw as a source of difference and is the more powerful design; it was
arrived at by accident, but the comparative claims in this work --- the effort
sweep, the ablations, the algorithm ranking --- rest on it and are unaffected.

What needed correcting is the \emph{absolute} values and their stated
uncertainty, not the comparisons between them.

\subsection{Computing Hardware}

All results in this thesis were produced on a single consumer laptop:

\begin{itemize}
  \item \textbf{CPU}: AMD Ryzen~7 4800H, 8 physical cores / 16 logical
        processors, \SI{2.9}{\giga\hertz} base.
  \item \textbf{RAM}: \SI{15.4}{\giga\byte}.
  \item \textbf{GPU}: none used. All training and inference ran on the CPU;
        the networks are small enough that this is not the binding constraint.
  \item \textbf{Thread budget}: capped at \textbf{8 of 16 logical processors}.
        Sustained all-core load caused the machine to stop responding entirely
        (\cref{sec:defects}), and benchmarking showed the additional parallelism
        yielded little throughput in any case --- eight threads are only
        $1.78\times$ faster than one for networks of this size.
  \item \textbf{Cumulative compute}: approximately 30~hours, comprising the
        hyperparameter search ($\approx 8$~h), the five-configuration benchmark
        ($\approx 6$~h), four pilot runs at $3 \times 10^5$ steps
        ($\approx 6$~h) and the reward and effort studies.
\end{itemize}

The modest hardware is directly responsible for the seed count, for the
$10^5$-step comparison budget, and for the decision not to tune every algorithm
individually. These constraints are stated here so that the scope of the
conclusions can be judged against them.

% -----------------------------------------------------------
\section{Defects Identified During Development}
\label{sec:defects}

Six substantive defects were identified in the model, the reward function and the
training configuration during the course of this work. They are documented here
rather than omitted, for two reasons: four of them produced results that appeared
entirely plausible while being incorrect, and the measurements that eventually
exposed them are part of the methodological contribution of this thesis.

\subsection{Model specified in incorrect angular units}

The \MuJoCo{} model format defaults to degrees unless declared otherwise, while
the joint ranges had been authored in radians. Every range was consequently
interpreted as under $3^{\circ}$ and the arm was effectively immobilised. The
symptom was a uniform 0\,\% success rate across all algorithms, which is readily
misread as a learning failure. All experiments predating the correction are
invalid.

\subsection{Compounding domain randomisation}

Body masses and muscle strengths are perturbed each episode to assess robustness.
The implementation multiplied the current values rather than rescaling from
nominal ones, producing a multiplicative random walk: after approximately one
thousand episodes, body mass had drifted to roughly 1\,\% of its correct value.
Caching nominal parameters and rescaling from them raised the measured drift
statistic from 0.013 to 0.998, where unity denotes no drift.

A second defect in the same routine invoked the random generator without a
\texttt{size} argument, returning a single scalar applied to every body and every
muscle. The intended independent per-element variation was therefore a single
global scale factor, which a policy can detect and compensate trivially, rendering
the robustness claim vacuous.

\subsection{Reward exploitable by episode termination}
\label{sec:exploit}

Under the initial reward formulation every per-step reward was non-positive, and
numerical divergence terminated the episode without penalty. In a discounted
formulation a terminated episode contributes no further return, so its value was
exactly zero, whereas surviving the full hundred steps accumulated approximately
$-20$. Terminating the episode by destroying the simulation was therefore the
optimal policy, worth roughly twice the success bonus and considerably easier to
attain.

Diagnostic measurement over fifty evaluation episodes established that the policy
had learned precisely this: 100\,\% of episodes ended in deliberate divergence at
approximately step 7 of 100, and none ran to completion.

The defect escaped detection because no recorded metric could expose it. The
reported error remained a plausible distance and continued to improve slowly, and
neither episode length nor divergence rate was logged --- so an agent abandoning
the task after 7 of 100 steps was indistinguishable from one attempting it and
plateauing. Both quantities are now recorded at every evaluation checkpoint.

The revised reward makes every per-step reward strictly positive, so early
termination forfeits value; pays the success bonus per step while the target is
held rather than once on contact, converting the objective to reach-and-hold; and
penalises divergence explicitly.

\subsection{Divergence masked by simulator self-recovery}

\MuJoCo{} resets its state internally upon detecting divergence. By the time
control returned to the environment the state was finite again, with the arm
returned to its rest configuration, so a test for non-finite values almost never
triggered. The episode continued, measuring the distance from the rest pose to the
target and recording it as data. Detection now uses the solver's warning
counters, which are the only reliable indicator, and the last finite error is
reported rather than the post-reset value.

\subsection{Replay ratio a factor of four below intended}
\label{sec:replaybug}

Training used four parallel environments with a gradient-step setting of one. In
the Stable-Baselines3 implementation, a single rollout iteration over a
four-environment vector collects four transitions and is followed by one gradient
update. Direct measurement confirmed the consequence:

\begin{table}[H]
\centering
\caption{Gradient updates per environment step, measured directly.}
\begin{tabular}{@{} L{8.0cm} C{4.0cm} @{}}
\toprule
\textbf{Configuration} & \textbf{Updates per step} \\
\midrule
As originally configured (4 envs, \texttt{gradient\_steps}=1) & 0.233 \\
Single environment, one step (standard) & 0.933 \\
4 envs, \texttt{gradient\_steps}$=4$ or $-1$ & 0.933 \\
\textbf{4 envs, \texttt{gradient\_steps}$=8$ --- the setting actually used} & \textbf{1.867} \\
\bottomrule
\end{tabular}
\label{tab:replayratio}
\end{table}

The last row is the one that matters for reading the results. The experiments do
not run at the ``restored'' ratio of $0.933$ but at $1.867$, because every run
script sets \texttt{gradient\_steps} $= 2 \times$ \texttt{N\_ENVS} explicitly
rather than relying on the module default. That is a deliberate choice --- a
replay ratio near two is common for off-policy control --- and it is applied
uniformly to every off-policy configuration, so it is not a confound between
them. It is recorded here because \cref{tab:hparams} reports
\texttt{gradient\_steps} $=8$ and the two statements must be read together.

Every run had therefore performed approximately 233\,000 gradient updates against
a nominal budget of 1\,000\,000 steps. An independent consistency check supports
this: 233\,000 updates at the measured \SI{17}{\milli\second} each predicts
71 minutes, against observed run times of 78.8 and 80.6 minutes.

The correction has a substantial cost. At the corrected ratio, one million steps
requires approximately 3.7 hours per seed, so a four-algorithm, ten-seed protocol
would require roughly 148 hours rather than the 50 previously assumed. That
protocol was never feasible on the available hardware; it appeared so only because
the runs were under-training.

\subsection{Default exploration temperature preventing convergence}
\label{sec:entropydefect}

Following the preceding five corrections the controller reached the target region
reliably but ceased improving, with closest approach remaining near
\SI{0.104}{\meter} across three distinct reward formulations and the fourfold
increase in effective training. The cause was the \SAC target entropy, whose
default of $-\dim(\mathcal{A})$ gives $-9$ for a nine-muscle action space.
Hyperparameter optimisation selected approximately $-19$, and the error halved.
The result is reported in \cref{sec:tuning}.

\subsection{Infrastructure failures}

Two non-scientific failures are recorded for completeness. A forced operating
system restart destroyed a 75-minute training run five minutes from completion,
because model state was written only on normal termination; training now
checkpoints at fixed intervals using atomic file replacement. Separately, the
machine ceased responding under four concurrent search processes, with no crash
dump written despite dumps being enabled and no hardware error logged, indicating
a thermal or power limit rather than a software fault. Subsequent work was
restricted to half the available logical processors. Benchmarking established
that the parallelism had provided little benefit in any case, thread scaling being
poor for networks of this size.

\section{Reproducibility and Sharing}
\label{sec:reproducibility}

Following the reproducibility guidelines of~\parencite{Pineau2021} for \RL{},
we ensure:
\begin{itemize}
  \item \textbf{Source code}: published under the MIT license on GitHub, with a
        pinned \texttt{requirements.txt} and the \texttt{git} hash recorded in each
        TensorBoard run.
  \item \textbf{Seeds and determinism}: all seeds are specified explicitly;
        \texttt{torch.use\_deterministic\_algorithms(True)} is enabled wherever
        possible.
  \item \textbf{Trained models}: all $40$ final policies
        ($4$~algorithms $\times$ $10$~seeds) are archived together with the
        associated \texttt{VecNormalize} statistics.
  \item \textbf{Evaluation logs}: the complete Optuna, TensorBoard, and CSV
        evaluation logs are provided.
  \item \textbf{Docker container}: a \texttt{Dockerfile} pinning the versions of
        \MuJoCo{}, PyTorch, and CUDA is provided.
\end{itemize}
